\begingroup
\small
\setlength{\tabcolsep}{3.5pt}
\renewcommand{\arraystretch}{1.15}
\begin{longtable}{@{}p{0.065\linewidth}p{0.285\linewidth}p{0.60\linewidth}@{}}
\caption{Normal-tier task catalogue. Method set A applies to every task in this table.}\label{tab:catalogue-normal}\\
\toprule
ID & Chemical task & Stated conditions and task-specific research target \\
\midrule
\endfirsthead
\multicolumn{3}{l}{\tablename\ \thetable\ (continued)}\\
\toprule
ID & Chemical task & Stated conditions and task-specific research target \\
\midrule
\endhead
\midrule
\multicolumn{3}{r}{Continued on the next page}\\
\endfoot
\bottomrule
\endlastfoot
NB1 & malonic acid thermal decarboxylation: malonic acid -> acetic acid + CO2\newline {\footnotesize Charge 0; multiplicity 1.} & \textit{Conditions:} gas-phase thermolysis.\newline \textit{Target:} Identify intermediates and preserve elementary-step continuity through decarboxylation to the final products. \\
\addlinespace[3pt]
NI1 & imine formation: formaldehyde + methylamine -> N-methylmethanimine + water\newline {\footnotesize Charge 0; multiplicity 1.} & \textit{Conditions:} gas phase; optimal around pH 4-6 (bell-shaped rate profile).\newline \textit{Target:} Identify the addition/condensation intermediates and connecting elementary steps leading to the imine and water. \\
\addlinespace[3pt]
NI5 & imine formation (aqueous): formaldehyde + methylamine -> N-methylmethanimine + water\newline {\footnotesize Charge 0; multiplicity 1.} & \textit{Conditions:} aqueous solution, weakly acidic (pH 4-6, bell-shaped rate profile).\newline \textit{Target:} Resolve the aqueous imine-formation network and relevant proton-transfer steps under the stated conditions. \\
\addlinespace[3pt]
NQ1 & aromatic Claisen rearrangement: allyl phenyl ether -> 2-allylphenol\newline {\footnotesize Charge 0; multiplicity 1.} & \textit{Conditions:} gas phase; thermal, \(\sim\)200 \({}^{\circ}\mathrm{C}\).\newline \textit{Target:} Resolve rearrangement and subsequent conversion to 2-allylphenol, retaining the intermediate network. \\
\addlinespace[3pt]
NQ5 & oxy-Cope rearrangement: 1,5-hexadien-3-ol -> hex-5-enal ([3,3] then enol-keto tautomerization)\newline {\footnotesize Charge 0; multiplicity 1.} & \textit{Conditions:} gas phase, heated.\newline \textit{Target:} Resolve the [3,3] rearrangement and enol--keto tautomerization as the investigated network. \\
\addlinespace[3pt]
NR1 & formamide dehydration: formamide -> hydrogen cyanide + water\newline {\footnotesize Charge 0; multiplicity 1.} & \textit{Conditions:} Gas phase; isolated neutral, closed-shell formamide; no added water, acid, or catalyst.\newline \textit{Target:} Investigate formimidic-acid intermediates and conformational transitions in the neutral dehydration network. \\
\addlinespace[3pt]
NU3 & H3P=CH2 ylide + formaldehyde -> ethylene + H3P=O\newline {\footnotesize Charge 0; multiplicity 1.} & \textit{Conditions:} salt-free model conditions (gas-phase ylide, closest to model chemistry).\newline \textit{Target:} Resolve the salt-free model Wittig-reaction network to ethylene and phosphine oxide. \\
\addlinespace[3pt]
NU5 & Passerini reaction: formaldehyde + formic acid + methyl isocyanide -> N-methyl-alpha-formyloxyacetamide\newline {\footnotesize Charge 0; multiplicity 1.} & \textit{Conditions:} dichloromethane (DCM), room temperature, high concentration.\newline \textit{Target:} Resolve the three-component Passerini-reaction network to the specified alpha-acyloxy amide. \\
\addlinespace[3pt]
NU6 & Staudinger beta-lactam synthesis: ketene + methanimine -> 2-azetidinone\newline {\footnotesize Charge 0; multiplicity 1.} & \textit{Conditions:} ketene generated in situ; inert solvent.\newline \textit{Target:} Identify the ketene--imine addition and beta-lactam-forming network, retaining any identified intermediate. \\
\addlinespace[3pt]
NY3 & isocyanic acid hydrolysis: HNCO + one water -> ammonia + carbon dioxide\newline {\footnotesize Charge 0; multiplicity 1.} & \textit{Conditions:} Gas phase; neutral closed-shell system with exactly one water molecule; no added solvent, acid, counterion, or catalyst.\newline \textit{Target:} Investigate carbamic-acid intermediates and connecting conformers without presupposing a fixed number of steps. \\
\addlinespace[3pt]
\end{longtable}
\endgroup
